\documentclass[11pt,a4paper]{article}
\usepackage[T1]{fontenc}
\usepackage[utf8]{inputenc}
\usepackage{lmodern}
\usepackage{amsmath,amssymb}
\usepackage[margin=25mm]{geometry}
\usepackage[hidelinks]{hyperref}
\hypersetup{pdftitle={Quantum preparations select continuations at a bounded Newton cusp},pdfauthor={navstokgap research working notes}}
\newcommand{\tightlist}{\setlength{\itemsep}{0pt}\setlength{\parskip}{0pt}}
\setlength{\emergencystretch}{3em}
\setcounter{secnumdepth}{0}
\begin{document}
\hypertarget{quantum-preparations-select-continuations-at-a-bounded-newton-cusp}{%
\section{Quantum preparations select continuations at a bounded Newton
cusp}\label{quantum-preparations-select-continuations-at-a-bounded-newton-cusp}}

\textbf{Result, 2026-10-27; written derivation, unrefereed.} For the
bounded \(C^1\) cusp potential already used in the
\href{newton-cusp-preparation-floor.md}{preparation-floor note}, there
are two explicit coherent-state families with the \textbf{same position
density} at \(t=0\), the same limiting phase-space point \((0,0)\), and
the same limiting incoming trajectory for \(t<0\). Their unique
Schrödinger evolutions have different outgoing semiclassical limits for
\(t>0\): one rebounds on the positive side and the other crosses to the
negative side. The preparations differ by a momentum shift tending to
zero. Their normalized superpositions yield any prescribed mixture of
the two outgoing point measures.

The result is local in time and uses the specified logarithmically
shrinking preparation. It neither classifies every possible preparation
nor selects a waiting-time distribution. It shows explicitly that the
limiting incoming phase-space history alone does not determine the
outgoing limit of these quantum families. The positive value of
\(\hbar\) is supplied, not derived.

This phenomenon has substantial prior art.
\href{https://arxiv.org/pdf/1011.1651}{Athanassoulis--Paul, Theorem 2.1
and Lemma 5.3} (theorem statements and relevant passages read) treat
splitting and one-dimensional cutoff \(-|q|^{1+\theta}\) potentials,
including \(\theta=1/2\), with detailed assumptions on phase-space
preparations. \href{https://arxiv.org/pdf/1608.03716}{Chabu, Theorem
1.12 and Proposition 3.1} (statements and surrounding discussion read)
exhibit the same incoming semiclassical history with distinct outgoing
histories, and mixtures, for the conical potential \(-|q|\). In Theorem
1.12 the two initial profiles can be chosen equal, so identical initial
position densities also occur in that conical example. The present
calculation uses the repository's smoother \(C^1\) cusp and explicit
Gaussian preparations; exact coverage by the cited rough-potential
theorems is not settled and no literature priority is claimed.

\hypertarget{potential-initial-data-and-target-trajectories}{%
\subsection{1. Potential, initial data and target
trajectories}\label{potential-initial-data-and-target-trajectories}}

Use dimensionless calibrated units, \(m=1\), and choose \(R>0\) and a
smooth even cutoff \(\chi\) equal to one on \([-R,R]\) and zero outside
\([-2R,2R]\). Put \[
 V(q)=-\frac23|q|^{3/2}\chi(q),\qquad
 F(q)=-V'(q),\qquad
 H_\hbar=-\frac{\hbar^2}{2}\partial_q^2+V(Q).\tag{1}
\] \(V\) and \(F\) are bounded, \(V\in C^1\), and on \(|q|\le R\),
\(F(q)=\operatorname{sgn}(q)\sqrt{|q|}\). The bounded-perturbation
theorem gives a self-adjoint \(H_\hbar\) on \(H^2(\mathbb R)\) for every
\(\hbar>0\). Write \(F_*=\|F\|_\infty\), and fix \(T>0\) so small that
\(F_*T^2/2<R/4\). All mean trajectories considered below remain in
\(|q|<R/2\) for small \(\hbar\).

There are constants \(C_0,C_1>0\), depending on the fixed cutoff, such
that for all \(x\) and all \(q\ne0\), \[
 |F(x)|\le C_0\sqrt{|x|},\qquad
 |F(x)-F(q)|
   \le C_1(1+|q|^{-1/2})\,|x-q|.\tag{2}
\] For \(|x-q|\le |q|/2\), the second bound follows from the derivative
of the square-root force and the smooth cutoff. Otherwise use the first
bound and \(|x-q|\ge |q|/2\). This is a Lipschitz bound anchored at
\(q\ne0\), \textbf{not} a global Lipschitz bound at the cusp.

For \(0<\hbar<e^{-1}\) set \[
 a_\hbar=\frac1{\log(1/\hbar)},\qquad
 g_{\hbar,k}(q)=(\pi\hbar)^{-1/4}
 \exp\!\left[-\frac{(q-a_\hbar)^2}{2\hbar}
     -\frac{i k a_\hbar^{3/4}(q-a_\hbar)}{\hbar}\right].\tag{3}
\] Choose a fixed \(K>0\) with
\[K^2>\max\{16\sqrt2\,C_0,\;16/3\}.\tag{4}\] The two preparations are
\(g_{\hbar,0}\) and \(g_{\hbar,K}\). Both have exactly the same Gaussian
position density and
\(\operatorname{sd}Q=\operatorname{sd}P=\sqrt{\hbar/2}\). Their means
are \[
 (q,p)=(a_\hbar,0),\qquad
 (q,p)=(a_\hbar,-K a_\hbar^{3/4}),\tag{5}
\] both converging to \((0,0)\). The classical energies of these
\textbf{centres} are respectively \(-\frac23a_\hbar^{3/2}\) and
\((K^2/2-2/3)a_\hbar^{3/2}\): both vanish, but their signs differ.

Let \(z_+(t)=(t^4/144,t^3/36)\) for \(-T\le t\le T\), and let \[
 z_\times(t)=
 \begin{cases}
  (t^4/144,t^3/36),&t\le0,\\
  (-t^4/144,-t^3/36),&t\ge0.
 \end{cases}\tag{6}
\] The first trajectory returns to the positive side, while the second
crosses to the negative side. Both follow the same incoming path, arrive
at \((0,0)\) at \(t=0\), and satisfy the cusp Newton equation.

\textbf{Theorem.} Let
\(\psi_{\hbar,k}(t)=e^{-itH_\hbar/\hbar}g_{\hbar,k}\). For \(k=0,K\),
the packet means and the position and momentum standard deviations
converge uniformly on \([-T,T]\): \[
 \begin{aligned}
 (\langle Q\rangle_{\hbar,0},\langle P\rangle_{\hbar,0})
       &\longrightarrow z_+(t),\\
 (\langle Q\rangle_{\hbar,K},\langle P\rangle_{\hbar,K})
       &\longrightarrow z_\times(t),\\
 \operatorname{sd}Q_{\hbar,k}(t)+\operatorname{sd}P_{\hbar,k}(t)
       &\longrightarrow0.
 \end{aligned}\tag{7}
\] In particular, the Weyl expectations of every
\(b\in C_c^\infty(\mathbb R^2)\) converge uniformly to \(b(z_+(t))\) and
\(b(z_\times(t))\), respectively.

\hypertarget{a-mean-and-spread-estimate-with-an-anchored-force}{%
\subsection{2. A mean and spread estimate with an anchored
force}\label{a-mean-and-spread-estimate-with-an-anchored-force}}

Suppress the subscripts, let \(q(t)=\langle Q\rangle\),
\(p(t)=\langle P\rangle\), and write
\(W(t)=\operatorname{sd}Q(t)+\operatorname{sd}P(t)\). The standard
Ehrenfest identities and Cauchy--Schwarz give \[
 \dot q=p,\qquad \dot p=\langle F(Q)\rangle,\qquad
 D^+W\le\operatorname{sd}P+\operatorname{sd}F(Q).\tag{8}
\] These identities hold as quadratic-form identities for the present
\(H_\hbar\): \(V\in W^{1,\infty}\), the Gaussian lies in
\(D(Q)\cap H^2\), and the required weighted and \(H^1\) domains are
preserved by the bounded-force evolution. They can be obtained by
regularizing \(Q,P\) and removing the cutoffs.

For \(q(t)\ne0\), (2) gives \[
 \operatorname{sd}F(Q)
 \le\|[F(Q)-F(q)]\psi\|
 \le L(q)\operatorname{sd}Q,\quad
 L(q):=C_1(1+|q|^{-1/2}),\tag{9}
\] and also \(|\langle F(Q)\rangle-F(q)|\le L(q)W\). Consequently \[
 D^+W\le[1+L(q)]W,\qquad W(0)=\sqrt{2\hbar}.\tag{10}
\] On an interval reached from the initial time without a cusp crossing,
if \(|q|\ge a_\hbar/2\) and its length is at most \(2T\), (10) gives \[
 W(t)\le\sqrt{2\hbar}\,e^{C_T/\sqrt{a_\hbar}},\qquad
 \int L(q(t))W(t)\,dt
    \le C_Ta_\hbar^{-1/2}\sqrt\hbar\,
                  e^{C_T/\sqrt{a_\hbar}}
    =o(a_\hbar^N)\tag{11}
\] for every fixed \(N>0\). Here and below \(C_T\) may change from line
to line but is independent of \(\hbar\). This super-polynomial smallness
in \(a_\hbar\) is why the logarithmic separation in (3) is convenient.
After the controlled passage in §3, its additional factor
\(e^{Ca_\hbar^{-1/4}}\) can be absorbed into \(e^{C_T/\sqrt{a_\hbar}}\),
so the same form of bound applies on the outgoing interval.

On the positive side, \(F(q)\ge0\) within the chosen core. For
\(g_{\hbar,0}\) at positive times, or either preparation at negative
times, (11) and (8) bootstrap \(q(t)\ge a_\hbar/2\): the total
force-bias error is \(o(a_\hbar)\), whereas the force itself pushes away
from the cusp in the appropriate time direction. On those intervals
compare \((q,p)\) with the classical solution from (5). The field is
Lipschitz on \(q\ge a_\hbar/2\) with constant \(O(a_\hbar^{-1/2})\).
Gronwall and (11) show that their phase-space distance is
\[O\!\left(a_\hbar^{-1/2}\sqrt\hbar\,
              e^{C_T/\sqrt{a_\hbar}}\right)=o(a_\hbar^N)
 \quad\text{for every fixed }N.\tag{12}\] The displayed constant \(C_T\)
absorbs the additional Gronwall factor.

\hypertarget{the-transmitted-packet-crosses-the-cusp}{%
\subsection{3. The transmitted packet crosses the
cusp}\label{the-transmitted-packet-crosses-the-cusp}}

It remains to control the \(k=K\) preparation while its mean passes
through \(q=0\), where (11) cannot be used with a fixed cutoff. Start at
\(q=a_\hbar\) and \(p=-Ka_\hbar^{3/4}\). Bootstrap, until the first
arrival at \(q=-a_\hbar\), the two inequalities
\[W\le a_\hbar,\qquad p\le-\frac K2a_\hbar^{3/4}.\tag{13}\] While
\(|q|\le a_\hbar\), the first inequality and (2) imply \[
 |\langle F(Q)\rangle|\le C_0\sqrt{\mathbb E|Q|}
 \le C_0\sqrt{2a_\hbar}.\tag{14}
\] The second inequality bounds the passage time by
\(4a_\hbar^{1/4}/K\). In that time the total change in \(p\) is at most
\((4\sqrt2 C_0/K)a_\hbar^{3/4}\), which is strictly below
\((K/2)a_\hbar^{3/4}\) by (4).

The same velocity bound makes the singular coefficient integrable: \[
 \int_{\rm passage}|q(t)|^{-1/2}\,dt
 \le \frac{2}{K a_\hbar^{3/4}}
        \int_{-a_\hbar}^{a_\hbar}|q|^{-1/2}\,dq
 =\frac8K a_\hbar^{-1/4}.\tag{15}
\] The mean crosses zero at most once, so (10) holds away from that
single time; apply it on the two sides and use continuity at the
crossing. It follows that
\[W\le\sqrt{2\hbar}\,e^{C a_\hbar^{-1/4}}=o(a_\hbar).\tag{16}\] The
inequalities (13) therefore close by continuity. The exit time
\(\tau_\hbar\) satisfies \(\tau_\hbar=O(a_\hbar^{1/4})\), and at exit
\[q(\tau_\hbar)=-a_\hbar,\qquad
 -\tfrac{3K}{2}a_\hbar^{3/4}
 \le p(\tau_\hbar)\le-\tfrac K2a_\hbar^{3/4}.\tag{17}\]

After this exit the force is negative in the chosen core. The same
bootstrap as in §2 keeps \(q\le-a_\hbar/2\) and makes (11)--(12) valid
against the classical solution started at the \textbf{actual} exit mean
(17). No self-adjoint Hamiltonian with a regularized potential is
substituted for (1).

\hypertarget{the-classical-limits-and-the-mixtures}{%
\subsection{4. The classical limits and the
mixtures}\label{the-classical-limits-and-the-mixtures}}

The zero-energy positive-side arrival solves
\(\dot q=-\sqrt{\frac43q^{3/2}}\). Integration gives \[
 |t|=\int_0^{q(t)}
 \frac{ds}{\sqrt{\frac43s^{3/2}}}
 =2\sqrt3\,q(t)^{1/4},\qquad
 q(t)=\frac{t^4}{144}.\tag{18}
\] For the \(k=0\) centre, energy conservation gives on either side of
\(t=0\) \[
 |t|=\int_{a_\hbar}^{q_a(t)}
 \frac{ds}{\sqrt{\frac43(s^{3/2}-a_\hbar^{3/2})}}.\tag{19}
\] The part from \(a_\hbar\) to \(2a_\hbar\) is \(O(a_\hbar^{1/4})\); on
the remainder the integrand is bounded by \(Cs^{-3/4}\) and tends to the
integrand in (18). Thus (19) tends to (18).

For \(k=K\) at negative times, and for its outgoing classical comparison
after (17), the analogous integral has denominator \[
 \sqrt{u_\hbar^2+\frac43(s^{3/2}-a_\hbar^{3/2})},
 \qquad u_\hbar\asymp a_\hbar^{3/4}.\tag{20}
\] Condition (4) ensures \(u_\hbar^2\ge(4/3)a_\hbar^{3/2}\) in the
outgoing case, so the integrand is dominated by \(Cs^{-3/4}\); it tends
to the one in (18). At the entrance for negative times, \(u_\hbar\) is
exactly \(Ka_\hbar^{3/4}\) and the same argument applies. The exit-time
shift tends to zero. Energy conservation gives the corresponding
momentum limits, and the bounds in §§2--3 supply (7), including a
uniform estimate on the shrinking passage interval.

The overlap of the two initial Gaussians is
\[\langle g_{\hbar,0},g_{\hbar,K}\rangle
 =\exp\!\left[-\frac{K^2a_\hbar^{3/2}}{4\hbar}\right]
 \longrightarrow0.\tag{21}\] It remains unchanged by unitary evolution.
For \(0\le\lambda\le1\) and fixed \(\theta\), normalize \[
 \Psi_{\hbar,\lambda,\theta}(0)
 \propto\sqrt\lambda\,g_{\hbar,0}
       +e^{i\theta}\sqrt{1-\lambda}\,g_{\hbar,K}.\tag{22}
\] For any \(b\in C_c^\infty(\mathbb R^2)\), the Fourier representation
of Weyl quantization and \(\|(e^{iA}-1)\psi\|\le\|A\psi\|\) show that \[
 \|[\operatorname{Op}^W_\hbar(b)-b(z_k(t))]
                         \psi_{\hbar,k}(t)\|\longrightarrow0
 \quad\text{uniformly for }|t|\le T,\tag{23}
\] using (7). Thus the cross term in a Weyl expectation equals
\(b(z_k(t))\langle\psi_{\hbar,0},\psi_{\hbar,K}\rangle+o(1)\) and
vanishes uniformly. The pure states (22) have the semiclassical limit \[
 \lambda\delta_{z_+(t)}+(1-\lambda)\delta_{z_\times(t)}.\tag{24}
\] For \(t\le0\) this is a single point on the common incoming path; for
\(t>0\) it has the prescribed two branch weights. The result does not
assert that the finite-\(\hbar\) Wigner function is nonnegative.

\hypertarget{consequence-for-state}{%
\subsection{5. Consequence for STATE}\label{consequence-for-state}}

The cusp's unique quantum evolution need not choose one classical
continuation from the limiting incoming phase-space history. Even
identical initial position densities admit distinct continuations
through vanishing momentum offsets, and coherent superpositions can
produce any mixture of the two demonstrated branches. These explicit
preparations supplement the earlier point-merger obstruction and align
with established rough-potential selection results; they do not derive a
physical preparation rule, a waiting-time selection, or a necessary
positive action constant.
\end{document}
